% Компилировать движком XeLaTeX (на Overleaf: Menu -> Compiler -> XeLaTeX).
\documentclass[11pt,a4paper]{article}

\usepackage{fontspec}
\setmainfont{DejaVu Serif}[Scale=0.92]
\setmonofont{DejaVu Sans Mono}[Scale=0.82]  % моноширинный шрифт с кириллицей (для \texttt и \url)
\usepackage{polyglossia}
\setmainlanguage{russian}
\setotherlanguage{english}
\usepackage[a4paper,margin=2.5cm]{geometry}
\usepackage{amsmath,amssymb}
\usepackage{booktabs}
\usepackage[font=footnotesize,labelfont=bf]{caption}  % уменьшенный шрифт подписей к рисункам и таблицам
\usepackage{tabularx}
\newcolumntype{L}{>{\raggedright\arraybackslash}X}
\usepackage{graphicx}
\usepackage{microtype}
\usepackage[numbers,sort&compress]{natbib}
\usepackage[hidelinks]{hyperref}
\usepackage{xurl}  % переносы длинных URL по строкам

% Параметры вёрстки: устраняют вылетающие за поле строки без «рыхлой» юстировки.
\graphicspath{{figures/}}
\tolerance=1400
\emergencystretch=2.5em
\hbadness=2200
\frenchspacing

\title{\textbf{Интегрированная модель мира для стратегического\\планирования и сценарного анализа}}
\author{Н.~Сотириади}
\date{\today}

\begin{document}
\maketitle

\begin{abstract}
\noindent
В работе описана интегрированная имитационная модель мира, в которой в едином
вычислительном контуре связаны экономика, климат и углеродный цикл, ущерб от
потепления, природные ресурсы, общество и политика, межгосударственные отношения,
культура, а также механизмы кризисов и финансов. Модель построена по принципу
системы взаимодействующих стран (57 страновых агентов и ряд
региональных агрегатов) и эволюционирует годовыми циклами: ежегодно каждый блок обновляет
своё состояние, и результаты передаются в следующий год и в смежные блоки. Цель
работы --- представить инструмент
стратегического планирования и сценарного анализа (проигрывания согласованных сценариев
развития, сопоставления альтернативных решений политики, ансамблевой оценки
неопределённости и обнаружения слабых сигналов приближающихся структурных сдвигов, в том
числе для управления системными рисками). Изложены проблематика и место модели
среди существующих подходов, её архитектура и лежащая в основе каждого блока
математика и эмпирическая база, процедура калибровки и проверки адекватности с
полученными метриками, анализ чувствительности и ансамблевая оценка неопределённости, явный перечень ограничений, а также типовые сценарии применения. Экономический и климатический блоки откалиброваны по историческим данным
за 1990--2023~годы и проходят ретроспективную проверку; блоки общества, политики и
геополитики проверяются по более скромному критерию --- превосходству над наивными
эталонными прогнозами.

\medskip
\noindent\textbf{Ключевые слова:} интегрированная модель, стратегическое планирование,
сценарный анализ, оценка неопределённости, управление риском, слабые сигналы.
\end{abstract}

\section{Проблематика}

Стратегическое планирование и сценарный анализ на длинных горизонтах --- от государственной
политики до корпоративной стратегии --- требуют инструмента, в котором климатические,
экономические и социально-политические процессы рассматривались бы совместно, а не порознь;
тот же запрос возникает и в управлении системными рисками. На практике аналитик вынужден
опираться на набор разнородных моделей: климатические модели описывают физику атмосферы, но почти не
содержат финансов и политики; интегрированные модели оценки (например, семейства DICE,
RICE, FUND) соединяют экономику и климат, но, как правило, лишены суверенных финансов,
межгосударственных отношений и социальной динамики \citep{nordhaus2017}; макроэкономические
модели подробно описывают денежно-кредитную сферу, но не учитывают природный капитал.
В результате взаимное влияние сфер --- например, путь от потепления к падению урожаев,
росту цен, социальному напряжению и далее к политической нестабильности и конфликту ---
оказывается за пределами возможностей любой отдельной модели.

Между тем именно сопряжение этих каналов определяет, как разворачиваются долгосрочные
сценарии развития и где накапливаются системные риски. Отсутствие связной модели, охватывающей
одновременно климат, экономику, общество и геополитику, --- это содержательный пробел,
ощутимый в стратегическом планировании и риск-менеджменте, где требуется не оценить отдельный
шок, а проследить его распространение по взаимосвязанным подсистемам и сопоставить
альтернативные траектории развития. Настоящая работа описывает модель, построенную для
заполнения этого пробела.
Подчеркнём с самого начала принципиальную установку: модель не претендует на роль
безошибочного предсказателя. Оценки многолетних прогнозов показывают, что даже
профессиональные прогнозы роста на горизонте от двух до пяти лет нередко уступают наивной
экстраполяции среднего \citep{celasun2021}. Поэтому реалистичная цель состоит не в точечном
предсказании, а в \emph{корректной согласованности модели с её собственной мерой
неопределённости} и в способности связно проигрывать сценарии.

Новизна работы --- не в отдельном блоке, а в охвате связей. Интегрированные модели оценки
(DICE, FUND) и крупные модельные комплексы (GCAM, MESSAGEix, REMIND) детально соединяют
экономику и климат, но не содержат суверенных финансов, межгосударственных отношений и
социально-политической динамики; классические системные модели мира (World3) лишены явной
экономики цен и климата современного уровня. Отличительная черта предлагаемой модели ---
доведение контура до общества, геополитики и дискретных кризисов и встраивание агентов на
основе языковых моделей, играющих за страны (таблица~\ref{tab:compare}, раздел~\ref{sec:repo});
количественное сопоставление этого преимущества с секторальными моделями вынесено в
приложение~\ref{app:benchmark}.

\begin{table}[t]
\centering
\caption{Сопоставление охвата с типовыми интегрированными и системными моделями (схематично;
$\sim$ --- ограниченная поддержка).}
\label{tab:compare}
\small
\begin{tabular}{@{}lccccc@{}}
\toprule
Возможность & DICE & GCAM & MESSAGEix & World3 & GIM \\
\midrule
Климат и углеродный цикл & да & да & да & нет & да \\
Экономика и цены & да & да & да & $\sim$ & да \\
Природные ресурсы & нет & да & да & да & да \\
Общество и политика & нет & нет & нет & $\sim$ & да \\
Геополитика и конфликты & нет & нет & нет & нет & да \\
Дискретные кризисы и финансы & нет & нет & $\sim$ & нет & да \\
Агенты на основе языковых моделей & нет & нет & нет & нет & да \\
\bottomrule
\end{tabular}
\end{table}

Сопоставление охвата (таблица~\ref{tab:compare}) полезно дополнить сравнением по
\emph{методологическим осям}: перечисленные модели решают разные задачи, поэтому прямое
ранжирование «лучше--хуже» неуместно, а различие удобнее проводить по парадигме, трактовке
неопределённости и предмету проверки. По \emph{парадигме решения} семейство DICE\,/\,RICE\,/\,FUND
и REMIND строится на оптимизации общественного благосостояния (поиск оптимальной траектории
политики), GCAM и MESSAGEix --- на детальном равновесии энергетических систем и систем землепользования, World3 --- на системной динамике мировых агрегатов, тогда как предлагаемая модель ---
рекурсивная имитация взаимодействующих стран с адаптивными ожиданиями. Отсюда и разделение
сильных сторон: модели оценки и энергетические комплексы \emph{глубже} в своей области
(оптимальная политика, отраслевая детализация энергетики), и предлагаемая модель их не заменяет.
Её отличие состоит в следующем: \emph{(i)}~широта сопряжения, доводящая
единый контур до общества, геополитики и дискретных кризисов и количественно проверенная
против секторальных моделей в контролируемом сопоставлении (приложение~\ref{app:benchmark});
\emph{(ii)}~неопределённость как основной объект --- ансамбли и сопоставление с историей встроены в рабочий процесс, а не
добавлены постфактум; \emph{(iii)}~проверяемость нетипичных блоков (ранжирование риска конфликта
по AUC, идентифицированный запаздывающий канал между экономикой и обществом), которых в
сопоставимых моделях попросту нет. Тем самым модель занимает нишу инструмента широкого
сценарного охвата и анализа распространения шоков, дополняя, а не вытесняя специализированные
модели в их областях.

\section{Модель и её архитектура}

\subsection{Общее устройство}

Модель представляет мир как совокупность стран, эволюционирующих по годам. Состояние
каждой страны включает экономические, ресурсные, социальные, политические и военные
переменные, а также двусторонние отношения с другими странами; над странами надстроены
глобальные переменные климата и углеродного цикла и мировые цены на ресурсы. Один шаг
модели --- это один год: блоки обновляются в фиксированном порядке, при котором выход
одного блока служит входом следующего, после чего результаты года передаются в
следующий год. Расчётное ядро детерминировано; неопределённость вводится через
ансамблевые прогоны со случайным выбором параметров из обоснованных диапазонов
(раздел~\ref{sec:calib}).

Ниже описаны блоки модели, объединённые в содержательные группы. Для каждого указано,
что он вычисляет, на какой математике и каких данных основан и что он даёт остальной
модели.

\subsection{Экономика}

Выпуск каждой страны определяется тремя факторами: накопленным капиталом, занятым
населением и потребляемой энергией. Производственная функция задана в форме вложенной
функции с постоянной эластичностью замещения (тип \emph{капитал--труд--энергия}),
нормированной так, что в базовой точке калибровки она совпадает с производственной
функцией Кобба--Дугласа \citep{arrow1961}:
\begin{equation}
Y = A \, \Theta \, K^{\alpha} L^{\beta} E^{\gamma},
\qquad \alpha = 0{,}30,\; \beta = 0{,}60,\; \gamma = 0{,}042,
\label{eq:prod}
\end{equation}
где $A$ --- совокупная производительность факторов, $\Theta$ --- технологический
уровень. Вложенная форма с эластичностью замещения капитала и энергии $\sigma_{KE}=0{,}4$
описывает, насколько трудно заместить один фактор другим: при удорожании энергии хозяйство
сдвигается к капиталу и снижает энергопотребление. Именно этот механизм позволяет цене на
углерод реалистично перестраивать экономику. Капитал накапливается по стандартному
правилу
\begin{equation}
K_{t+1} = (1-\delta)\,K_t + s\,Y_t, \qquad \delta = 0{,}05,
\label{eq:capital}
\end{equation}
где норма сбережения $s\approx0{,}24$ модулируется устойчивостью и напряжённостью в
обществе. Производительность $A$ растёт с базовым темпом около одного процента в год,
дополнительно реагируя на расходы на исследования, переток технологий через торговлю и
член условной конвергенции, в котором более бедные экономики догоняют быстрее (наклон $0{,}0093$ на
лог-единицу разрыва душевого выпуска до границы, оценён по срезу реального ВВП по ППС за 2015--2023,
так что модель воспроизводит наблюдаемый догоняющий рост Китая и Индии); для прогнозов вперёд её
дрейф привязан к сценарию социально-экономического развития SSP2 \citep{riahi2017}. Рынки энергии,
ресурсов и капитала замыкаются через цены: спрос на энергию и инвестиции реагирует на
цену энергии и стоимость капитала с откалиброванной эластичностью. Отдельно
отслеживаются инфляция (кривая Филлипса с энергетическим переносом издержек),
безработица (закон Оукена), ставка центрального банка (правило Тейлора, \citep{taylor1993}),
бюджет и долг государства, а также торговые потоки. Дополнительно инфляция несёт слабый член
количественной теории денег: избыточный рост широкой денежной массы (эндогенного запаса банковских
депозитов) сверх эталона при стабильной скорости обращения переносится в цены, причём коэффициент
переноса ($\lambda\approx0{,}027$), оценённый по межстрановой панели, мал по построению --- что
согласуется с пологой после-1990 кривой Филлипса, --- а сама оценка независимо подтверждает
заложенный в модели якорь инерции инфляции. Экономический блок воспроизводит
последнее десятилетие реальной истории и потому является не умозрительным, а
эмпирически закреплённым.\footnote{Источники калибровки экономического блока: Penn World
Table (доли факторов производства и совокупная производительность); база World Development
Indicators Всемирного банка (валовой продукт, население, энергопотребление, военные расходы);
прогнозный дрейф производительности привязан к сценариям социально-экономического развития
SSP2 (база IIASA). Полный перечень переменных с версиями наборов и ссылками приведён в файле
\texttt{data/external/SOURCES.md} репозитория (раздел~\ref{sec:repo}).}

\subsection{Климат и углеродный цикл}

Выпуск сопровождается выбросами углекислого газа; к ним добавляется источник, связанный
с землепользованием. Углеродоёмкость выпуска структурно снижается с темпом, зависящим от страны и
растущим с уровнем развития: более богатые постиндустриальные экономики декарбонизируются быстрее
(возобновляемая энергетика и вынос тяжёлой промышленности) --- зеркальное отражение конвергенции
роста, оценённое по спаду углеродоёмкости по странам за 2015--2023. Выбросы накапливаются в атмосфере
и распадаются на нескольких характерных временных горизонтах, учитываются прочие парниковые газы, а
результирующее потепление рассчитывается стандартной двухрезервуарной моделью теплового баланса
поверхности и
океана. Блок откалиброван по общепринятой научной оценке: он воспроизводит принятую
чувствительность климата к удвоению концентрации углекислого газа (равновесная
чувствительность~$3{,}0$\,$^\circ$C) и отслеживает наблюдаемый ход температуры и
углерода с 1990 по 2023~год \citep{ipcc_ar6_wg1,friedlingstein2023}.\footnote{Источники
климатического блока: оценка равновесной чувствительности --- Шестой оценочный доклад МГЭИК
(рабочая группа~I); ряд глобальных выбросов --- Global Carbon Project (Global Carbon Budget) и
свод Our World in Data; ряд приземной температуры за 1990--2023~годы --- NOAA\,/\,HadCRUT5.}
Наиболее сильные
самоусиливающиеся обратные связи (таяние вечной мерзлоты, резкий выброс углерода) глубоко
неопределённы и поэтому по умолчанию исключены из основного прогона и исследуются лишь как
хвостовые риски. Это быстрая откалиброванная аппроксимация, уместная в модели экономики и
общества, но не заменяющая полноценную модель земной системы.

\subsection{Ущерб от потепления и стоимость углерода}

Потепление возвращается в экономику как ущерб: повышение температуры снижает выпуск.
Принята квадратичная функция ущерба, снижающая уровень выпуска,
\begin{equation}
\Omega(T) = 1 - 0{,}006\,T^{2},
\label{eq:damage}
\end{equation}
что соответствует потерям $2{,}4\,\%$, $5{,}4\,\%$ и $9{,}6\,\%$ валового продукта при
потеплении на $+2$, $+3$ и $+4\,^\circ$C. Эта оценка лежит в середине эмпирического
интервала: примерно в два с половиной раза выше, чем в DICE-2016R2, и ниже центральной
оценки мета-анализа Говарда и Стернера \citep{howard2017,burke2015}. Прямая межмодельная проверка
показывает, откуда берётся наша итоговая стоимость углерода: если собственному движку предельного
импульса задать \emph{собственный} заниженный коэффициент ущерба DICE и его дисконтирование, он
возвращает около \$20 за тонну на многовековом горизонте --- ниже \$31 у DICE-2016R2. Этот разрыв ---
не артефакт движка, а эффект роста: наше эмпирически калиброванное ядро роста (с конвергенцией по
уровню развития, раздел~\ref{sec:calib}) повышает душевое потребление быстрее экзогенной траектории
DICE, поэтому более богатое будущее сильнее дисконтирует свои ущербы. Более высокая итоговая
стоимость углерода следует из нашей более высокой, эмпирически обоснованной функции ущерба (примерно
в два с половиной раза выше DICE), а оценка ниже DICE на его же ущербах указывает, что DICE-2016R2
занижает ущербы, если задать реалистичный рост. На основе функции
ущерба вычисляется социальная стоимость углерода --- экономический ущерб от выброса
дополнительной тонны углекислого газа, стандартный ориентир климатической политики, ---
методом предельного импульса. Эта величина существенно зависит от схемы дисконтирования
и структуры будущего роста; результаты приведены в разделе~\ref{sec:calib}.\footnote{Калибровка
функции ущерба опирается на мета-анализ Говарда и Стернера и панельные оценки Бёрка с
соавторами; ориентиры социальной стоимости углерода --- оценки Агентства по охране окружающей
среды США (EPA, 2023) и проекта RFF-SP.}

\subsection{Ресурсы}

Для каждой страны отслеживаются запасы, добыча и потребление энергии, продовольствия и
металлов; мировые цены растут, когда спрос превышает предложение, и эта цена передаётся
в инфляцию и далее в экономику. Ресурсный рынок замыкается по цене с откалиброванной
эластичностью спроса; в качестве запасного предусмотрено правило медленной подстройки.
Ресурсы представлены укрупнённо, а не как детальные отраслевые энергосистемы.\footnote{Ресурсный
блок задан укрупнённо: энергетические показатели опираются на базу World Development Indicators;
ценовые эластичности заданы обобщёнными рыночными оценками, а не отдельным отраслевым набором
данных.}

\subsection{Общество, политика, геополитика и культура}

Для каждой страны модель отслеживает доверие к власти, социальное напряжение, неравенство
доходов, легитимность институтов и протестное давление. Эти переменные реагируют на
экономику --- рост инфляции и безработицы повышает напряжение и подрывает доверие, а
высокое неравенство усиливает этот эффект --- и, в свою очередь, влияют на устойчивость и
риск кризисов. Каждая пара стран связана отношением: объёмом торговли, наличием союза,
уровнем конфликта или напряжённости, состоянием войны и санкциями. Военный потенциал
каждой страны привязан к наблюдаемым показателям национальной мощи, а не к произвольно
назначенным значениям \citep{singer1972}. Небольшой набор культурных признаков (прежде всего степень
индивидуализма, избегание неопределённости, долгосрочность ориентации и иерархичность),
а также характер политического режима влияют на экономическое и социальное поведение
\citep{hofstede2010}. В отличие от экономики и климата, для этих связей нет единого
авторитетного калибровочного набора; там, где это возможно, их приоры привязаны к литературе и
воспроизводимым числам (раздел~\ref{sec:uq}), а остаток задан экспертно и проверяется по более
скромному критерию (раздел~\ref{sec:calib}).\footnote{Социально-политические
и геополитические переменные привязаны к признанным индексам: неравенство доходов --- база SWIID
и индекс Джини Всемирного банка; легитимность, качество управления и политическая стабильность ---
Worldwide Governance Indicators; вооружённые конфликты --- UCDP/PRIO Armed Conflict Dataset
(вер.~24.1); военный потенциал --- National Material Capabilities (Correlates of War, индекс CINC) и
база военных расходов SIPRI; санкции --- Global Sanctions Database (вер.~3); культурные признаки ---
индексы Хофстеде. Сопоставление каждой переменной с источником описано в \texttt{data/external/SOURCES.md}.}
Признаки, которые на деле не влияли на исходы, из модели удалены, чтобы она не
имитировала культурную детализацию, которую не использует.

\subsection{Риск, кризисы и финансы}

Поверх плавной годовой динамики модель распознаёт и применяет дискретные кризисы:
долговые, валютные, крах политического режима и климатические экстремальные события, ---
каждый из которых запускается при пересечении соответствующего порога. Тяжесть кризисов
может задаваться распределением с тяжёлым хвостом (большинство событий умеренны, редкие ---
катастрофичны), что отвечает наблюдаемому распределению реальных кризисов и войн
\citep{taleb2007,weitzman2009}; этот режим необязателен и исследуется как
неопределённость. Финансовая сторона ведётся со строгим учётом: всякое изменение долга
уравновешено встречным потоком. Государственные финансы замкнуты, а частный сектор
описан полным банковским балансом, в котором каждый кредит порождает встречный депозит,
так что денежная масса тождественно равна депозитам и кредитам; избыточное заимствование
повышает стоимость кредита и сдерживает экономику --- финансовый акселератор
\citep{bernanke1999,godley2007}.\footnote{Для блока риска и финансов отдельного калибровочного
набора данных нет: распределения тяжести кризисов с тяжёлым хвостом и балансовый (stock-flow
consistent) учёт опираются на методологическую литературу, указанную по тексту, а не на единый
внешний источник.}

\subsection{Взаимодействие блоков}

Связность --- определяющее свойство модели: 
экономика порождает выбросы; 
выбросы определяют потепление; 
потепление через функцию ущерба~\eqref{eq:damage} снижает выпуск;
цена на углерод удорожает энергию и через замещение в~\eqref{eq:prod} снижает выбросы;
экономические трудности повышают социальное напряжение и снижают доверие; 
это меняет норму сбережения в~\eqref{eq:capital} и повышает риск кризиса; 
кризис бьёт по выпуску и капиталу и передаётся партнёрам по торговле; 
ресурсный дефицит поднимает цены и инфляцию;
конфликты и санкции нарушают торговлю. 

Тем самым шок в одной подсистеме распространяется по остальным --- именно это сопряжение и составляет содержательную ценность модели для анализа рисков. Эта межподсистемная сеть географически заземлена: двусторонняя торговля подчиняется гравитационному закону (интенсивность убывает с расстоянием, эластичность ${\approx}0{,}9$; \citealp{headmayer2014}), а конфликты, социальная напряжённость и климатический риск перетекают к приграничным соседям, так что шок распространяется по реалистичным пространственным путям, а не равномерно. Включение этого географического сопряжения практически не меняет ретроспективное соответствие 2015--2023 (RMSE по продукту, $\mathrm{CO_2}$ и температуре сдвигаются менее чем на четверть процента), так что выигрыш в реализме распространения не стоит ничего в проверенной точности.

\begin{figure}[t]
\centering
\includegraphics[width=\linewidth]{fig1_architecture.pdf}
\caption{Архитектура модели и связи между блоками. Экономика выступает центральным
узлом; чёрные стрелки показывают прямые воздействия, красные --- обратные связи,
снижающие выпуск или сбережения. Замкнутый климатический контур (от выбросов к потеплению,
далее к ущербу и снижению выпуска) дополнен ценовым каналом углерода, ресурсным, социально-политическим
и геополитическим блоками, а также блоком дискретных кризисов. Нижний слой реализует
непрерывный мониторинг совместного состояния для выявления слабых сигналов.}
\label{fig:arch}
\end{figure}

\section{Калибровка и проверка адекватности}
\label{sec:calib}

\subsection{Источники данных}

Каждый блок привязан к внешним данным (таблица~\ref{tab:data}). Экономический и
климатический блоки откалиброваны и проверены наиболее строго; социально-политические и
геополитические переменные привязаны к признанным индексам, но проверяются по более
мягкому критерию.

\begin{table}[t]
\centering
\caption{Эмпирическая база блоков модели.}
\label{tab:data}
\small
\begin{tabularx}{\linewidth}{@{}l L L@{}}
\toprule
Блок & Источник данных & Назначение \\
\midrule
Экономика & Penn World Table, World Development Indicators & доли факторов, рост, долг \\
Климат и углерод & оценка МГЭИК AR6, Global Carbon Budget & чувствительность, ход $T$ и $\mathrm{CO_2}$ \\
Стоимость углерода & EPA-2023, RFF-SP & ориентир для сверки \\
Прогнозный рост & сценарии SSP2 & дрейф производительности вперёд \\
Конфликты & UCDP Georeferenced Event Dataset & проверка риска конфликта \\
Неравенство & SWIID & неравенство доходов \\
Институты & Worldwide Governance Indicators & легитимность, качество управления \\
Военная мощь & составной индекс национальной мощи & потенциал стран \\
Культура & индексы Хофстеде & культурные признаки \\
\bottomrule
\end{tabularx}
\end{table}

\subsection{Принципы калибровки и контроль целостности}

Калибровка следует нескольким правилам. Во-первых, каждый необязательный механизм
выполнен как переключатель, по умолчанию выключенный, если он не закреплён исторической
привязкой; так основное («штатное») поведение модели остаётся привязанным к данным, а
спекулятивные механизмы исследуются отдельно как неопределённость. Во-вторых, для оценки
параметров применяется не точечная подгонка, а \emph{сопоставление с историей}: вместо
поиска единственного «наилучшего» набора параметров отбрасываются области параметрического
пространства, несовместимые с наблюдениями, и остаётся апостериорная область «пока не
отвергнутых» значений --- стандартный подход для ресурсоёмких имитационных моделей
\citep{williamson2013}.
В-третьих, климатические коэффициенты, выводимые из исходного состояния, жёстко связаны с
контрольной суммой исходных данных и не могут быть изменены вручную; целостность этой связи
проверяется автоматически. Наконец, действует регрессионный контроль: значимые величины
ретроспективной проверки зафиксированы, и любое изменение ядра должно либо сохранять их с
точностью до допуска, либо сопровождаться явной перекалибровкой.

\subsection{Полученные метрики}

Сводка валидации по всем блокам модели приведена в таблице~\ref{tab:backtest}, а ключевые
результаты разобраны ниже. Ретроспективная проверка за 2015--2023~годы даёт средне\-квадратичные
ошибки по валовому продукту, мировым выбросам и температуре. После исправления инициализации
капитала и Кобб--Дугласово, и замкнутое ядро с вложенной функцией замещения достигают ошибки по
продукту около $0{,}62$ триллиона долларов; выигрыш от вложенной функции замещения с рыночным
замыканием теперь проявляется в выбросах (рис.~\ref{fig:valid}a), где он снижает ошибку с $1{,}87$
до $0{,}93$ гигатонны.

\begin{table}[t]
\centering
\caption{Сводка валидации по блокам модели: источник данных, критерий проверки и достигнутый
результат. Экономический и климатический блоки проверяются строгой ретроспективной проверкой;
для остальных критерий мягче --- сверка с опубликованными оценками, превосходство над наивным
эталоном или идентификация канала. Сравнение текущего ядра с прежним (Кобб--Дуглас) показано на
рис.~\ref{fig:valid}a.}
\label{tab:backtest}
\small
\begin{tabularx}{\linewidth}{@{}l l L L@{}}
\toprule
Блок & Источник & Критерий & Результат \\
\midrule
Экономика & PWT, WDI & ретро-RMSE 2015--2023 & ВВП $0{,}62$ трлн долл.\ (Кобб--Дуглас $0{,}62$) \\
Климат и углерод & МГЭИК AR6, GCB & ретро-RMSE 2015--2023 & $\mathrm{CO_2}$ $0{,}93$ Гт; $T$ $0{,}135\,^\circ$C \\
Стоимость углерода & EPA-2023, RFF-SP & сверка с публикуемыми; дисконтирование по Рамсею & ${\approx}\$95$/т$\mathrm{CO_2}$ (диапазон $95$--$280$); на заниженных ущербах DICE оценка ${\approx}\$20$, т.е.\ DICE-2016R2 занижает ущербы относительно нашей эмпирически калиброванной функции \\
Конфликты & UCDP/PRIO & ранжирование vs реестр & AUC $0{,}736$ $[0{,}59;\,0{,}86]$; BSS $+0{,}143$; $p{\approx}0{,}001$ \\
Социальный блок & SWIID, WGI & идентификация канала (разд.~\ref{sec:uq}) & значимый запаздывающий канал; регрессия $p<10^{-12}$ \\
\bottomrule
\end{tabularx}
\end{table}

Сопоставление с эталонными прогнозами: модель
\emph{превосходит} инерционный прогноз по температуре (оценка успешности около $+0{,}14$)
и незначительно --- по продукту, но \emph{уступает} наивной линейной экстраполяции по
продукту и выбросам на этом коротком и гладком интервале. Это закономерно и согласуется с
оценками профессиональных прогнозов \citep{celasun2021}: структурная модель повышает
успешность там, где работает динамика (климат), и проигрывает тривиальной экстраполяции
там, где краткосрочный гладкий агрегат трудно превзойти.

Для блока конфликта применён иной критерий. Показатель конфликтной предрасположенности
стран сопоставлен с историческим реестром вооружённых конфликтов за 1990--2023~годы по
57~странам \citep{sundberg2013}. Существенна сама постановка проверки: оценивается не
подогнанный к этим же данным выход, а \emph{входной} показатель риска, который вычисляется из
наблюдаемых характеристик страны и не настраивался по реестру конфликтов. Поэтому здесь нет
ни обучаемых на целевой переменной параметров, ни связанной с ними утечки информации, а
скользящее разбиение на обучение и контроль излишне; оценивается всё множество из 57~стран
целиком. Модель уверенно ставит реально затронутые конфликтом страны выше спокойных: площадь
под характеристической кривой составляет $0{,}736$ (бутстреп-интервал $95\,\%$: $[0{,}59;\,0{,}86]$,
$20\,000$ ресэмплов выборкой пар с возвращением \citep{efron1993}), а перестановочная проверка
значимости (случайная перестановка меток, $50\,000$ повторений) даёт $p\approx0{,}001$, то есть
ранжирование заведомо лучше случайного.
Коэффициент успешности по Брайеру относительно прогноза базовой частоты равен $+0{,}143$
(положительная величина означает наличие прогностической ценности). Это разумная планка
для блока такого рода: модель пригодна для оценки относительного риска, но не для
предсказания конкретных войн в конкретные даты.

\begin{figure}[t]
\centering
\includegraphics[width=\linewidth]{fig2_validation.pdf}
\caption{Результаты ретроспективной проверки. \emph{(a)}~Среднеквадратичная ошибка
ретропроверки 2015--2023 для замкнутого ядра (вложенная CES) в сравнении с прежним ядром
Кобба--Дугласа по мировому продукту, выбросам $\mathrm{CO_2}$ и температуре (меньше --- лучше).
\emph{(b)}~ROC-кривая конфликтного блока: ранжирование 57~стран по предрасположенности к конфликту
против реестра вооружённых конфликтов 1990--2023~годов даёт площадь под кривой $0{,}736$ (бутстреп
$90\,\%$ полоса показана заливкой; $95\,\%$ ДИ $[0{,}59;\,0{,}86]$; перестановочный тест
$p\approx0{,}001$) против $0{,}5$ у случайного угадывания.}
\label{fig:valid}
\end{figure}

Социальная стоимость углерода чувствительна к схеме дисконтирования. При современной
схеме Рамсея с околонулевой нормой временного предпочтения она составляет около
$95$ долларов за тонну --- в пределах широкого современного диапазона, ниже
центральных оценок EPA-2023 и RFF-SP (около $190$) \citep{rennert2022,epa2023}; при
дисконтировании по Нордхаусу --- около $42$ долларов. Существенно, что эта величина
\emph{действительно} чувствительна к структуре будущего роста: при разных вариантах
калибровки экономического ядра она меняется в диапазоне примерно $95$--$280$ долларов, а
рекалибровка ядра роста под конвергенцию по уровню развития (более быстрый догоняющий рост, который
делает будущее богаче и сильнее дисконтирует его ущербы) сдвинула центральное значение вниз примерно
на четверть. На собственном заниженном коэффициенте ущерба DICE-2016R2 наш движок при этом быстром
росте даёт около $20$ долларов --- ниже $31$ у DICE, --- так что DICE занижает ущербы относительно
принятой нами эмпирически калиброванной функции. Мы приводим эти значения вместе с их чувствительностью, а не подгоняем под
целевой ориентир; сама эта чувствительность представляет самостоятельный результат.

\section{Неопределённость, чувствительность и иденти\-фикация параметров}
\label{sec:uq}

Важный вопрос --- идентифицированы ли многочисленные параметры модели данными, а не заданы
вручную. Каждый ключевой параметр снабжён априорным распределением ---
центральным значением, диапазоном правдоподобия и литературным источником
(приложение~\ref{app:params}); оценка ведётся не точечной подгонкой, а сопоставлением с
историей (раздел~\ref{sec:calib}), при котором отбрасываются области параметрического
пространства, несовместимые с наблюдениями. Так, эластичность замещения капитала и энергии
$\sigma_{KE}=0{,}4$ опирается на эмпирические оценки в диапазоне $0{,}3$--$0{,}5$, норма
амортизации $\delta=0{,}05$ --- на данные Penn World Table ($3$--$8\,\%$), а коэффициент
ущерба $0{,}006$ лежит в середине эмпирического интервала $0{,}0015$--$0{,}025$, между оценкой
DICE-2016R2 и центральной оценкой Говарда и Стернера. Неустранимая неопределённость этих
величин не скрывается, а переносится в расчёт через ансамбли.

\textbf{Анализ чувствительности.} Чтобы установить, какие параметры действительно определяют
исходы, проведён глобальный скрининг методом элементарных эффектов Морриса
\citep{morris1991,campolongo2007} по 33~ключевым параметрам для четырёх выходных
величин (рис.~\ref{fig:sens}; $8$~траекторий, горизонт 10~лет). Картина согласуется с
устройством модели и физикой: температуру определяют почти исключительно климатические
параметры --- равновесная чувствительность, обмен теплом с океаном и теплоёмкость
поверхности; мировой продукт --- амортизация капитала, коэффициент ущерба и энергоёмкость
выпуска; мировые выбросы --- масштаб эмиссий и темп структурной декарбонизации. Эти выводы
устойчивы к бюджету скрининга: при увеличении числа траекторий с $8$ до $32$ ранговая корреляция
Спирмена полного вектора $\mu^*$ с эталонным прогоном составляет $\rho\approx0{,}99$ для трёх
физических величин (в диапазоне $0{,}986$--$0{,}995$), а доминирующие параметры воспроизводятся при
всех проверенных стартовых значениях генератора случайных чисел: для температуры первая пятёрка
совпадает полностью, а для выбросов и продукта устойчиво воспроизводится доминирующая тройка
(масштаб эмиссий, темп декарбонизации и землепользование --- для выбросов; амортизация капитала,
коэффициент ущерба и энергоёмкость выпуска --- для продукта), тогда как нижестоящие параметры
близки по величине и могут меняться местами. Социальное напряжение, напротив, к этим физико-экономическим параметрам почти
нечувствительно: его элементарные эффекты лежат у шумового порога скрининга (абсолютная мера
$\mu^*$ примерно на три порядка ниже, чем у физических величин), а само их
ранжирование, в отличие от устойчивого ранжирования физических величин, неустойчиво к выбору
стартового значения ($\rho\approx0{,}93$ при $8\!\to\!32$ траекториях против $\approx0{,}99$ у
физических величин). Важно, однако, что это \emph{не} означает отвязанности общества от
экономики: данный скрининг варьирует лишь физико-экономические приоры, тогда как напряжение
управляется собственными (социальными) параметрами и реагирует на экономические условия с
многолетним лагом. Этот канал влияния экономических условий на социальное напряжение мы
выделяем отдельно ниже (рис.~\ref{fig:social}). Тем самым
скрининг подтверждает, что значимые величины управляются содержательно уместными и эмпирически
закреплёнными параметрами, и этот вывод не является артефактом ограниченного бюджета расчётов.

\textbf{Ансамблевая оценка неопределённости.} Поскольку ядро детерминировано, неопределённость
вводится прогоном Монте-Карло ансамбля ($500$~членов) со случайным выбором параметров из их
априорных диапазонов. Результат --- не одна траектория, а распределение исходов
(рис.~\ref{fig:ens}). Здесь и далее «мировой продукт» означает \emph{совокупный валовой продукт
$57$~моделируемых акторов} --- $50$~крупнейших экономик, остаточные региональные агрегаты (все
прочие страны каждого региона) и Тайвань, --- которые вместе покрывают практически весь мировой
выпуск; в базовом $2023$~году (валидированный канон) он составляет около
$107$~трлн долларов. На горизонте десяти лет его медиана составляет около
$127$~трлн долларов при межквартильном размахе примерно $125$--$130$ и интервале $5$--$95$-го
перцентиля около $121$--$133$; интервал $5$--$95$-го перцентиля аномалии температуры составляет
примерно $1{,}0$--$2{,}0\,^\circ$C и отражает прежде всего разброс равновесной чувствительности.
Приведённые перцентили сошлись по объёму ансамбля: при росте числа членов с $80$ до $500$ медиана
и верхние перцентили продукта и температуры устойчивы, а Монте-Карло погрешность каждого
перцентиля (бутстреп по членам ансамбля) не превышает $\pm0{,}35$~трлн долларов для продукта и
$\pm0{,}03\,^\circ$C для температуры; единственная заметно зависящая от объёма выборки величина ---
нижний хвост температуры, уточняющийся с $\approx0{,}9$ до $\approx1{,}0\,^\circ$C. Именно
такие интервалы, а не точечные значения, и составляют содержательный результат модели.

\begin{figure}[t]
\centering
\includegraphics[width=\linewidth]{fig3_sensitivity.pdf}
\caption{Глобальный скрининг чувствительности (метод Морриса) по 33~ключевым параметрам для
четырёх выходных величин ($57$~стран); показаны семь наиболее влиятельных параметров, мера $\mu^*$
нормирована на максимум в каждой панели, а абсолютный максимум $\mu^*$ указан в рамке. Температура
управляется климатическими параметрами, продукт и выбросы --- экономическими и эмиссионными. У
социального напряжения абсолютный $\mu^*$ примерно на три порядка ниже (у шумового порога): к
физико-экономическим приорам оно почти нечувствительно, а его собственный канал разобран на
рис.~\ref{fig:social}.}
\label{fig:sens}
\end{figure}

\begin{figure}[t]
\centering
\includegraphics[width=\linewidth]{fig4_ensemble.pdf}
\caption{Ансамблевая оценка неопределённости ($500$~членов): распределение исходов по Монте-Карло
прогону со случайным выбором параметров из априорных диапазонов. Показаны медиана и полосы
$25$--$75$-го и $5$--$95$-го перцентилей для \emph{(a)}~совокупного продукта $57$~моделируемых
стран и \emph{(b)}~аномалии температуры.}
\label{fig:ens}
\end{figure}

\textbf{Канал влияния экономики на социальное напряжение.} Малая чувствительность напряжения к
физико-экономическим \emph{приорам} (выше) не означает, что общество отвязано от экономики:
напряжение --- накапливающаяся переменная (приращение определяется безработицей, инфляцией и
неравенством), поэтому связь с экономикой проявляется не мгновенно, а с лагом, и недоступна
скринингу фиксированного горизонта. Мы выделяем этот канал тремя взаимодополняющими проверками
(рис.~\ref{fig:social}; воспроизводимый расчёт описан в разделе~\ref{sec:repo}). Во-первых,
\emph{локализация}: в ансамбле, дополнительно варьирующем приоры самого социального блока, итоговое
напряжение значимо управляется именно ими --- чувствительностью к неравенству (ранговая корреляция
Спирмена $\rho\approx+0{,}43$, $p<10^{-6}$) и силой якоря доверия ($\rho\approx-0{,}50$,
$p<10^{-8}$), --- тогда как физико-экономические приоры остаются слабыми; то есть в рассмотренном
выше скрининге просто не варьировались прямые драйверы блока. Во-вторых, \emph{накопление}: вызванный приорами
разброс напряжения растёт с горизонтом (межчленное стандартное отклонение увеличивается примерно
в полтора раза от $5$ к $25$~годам), что и есть подпись лага. В-третьих, \emph{импульсный отклик}:
экзогенный стагфляционный шок (рост безработицы и инфляции на три года) повышает напряжение, причём
отклик нарастает в течение первых $\approx4$~лет (положителен у $93$--$100\,\%$ прогонов ансамбля
приоров) и сохраняется более десятилетия; на дальних горизонтах разброс растёт, так как у части
прогонов шок запускает дискретные кризисы. Распределённо-лаговая регрессия приращения напряжения на
экономический стресс (внутригрупповые фиксированные эффекты, кластерные ошибки) подтверждает
значимую передачу на лагах $0$--$1$~год ($p<10^{-12}$). Таким образом, канал влияния экономики на
напряжение статистически значим, но мультиплицирован и запаздывает; именно это запаздывание, а не
отсутствие связи, объясняет нулевой результат скрининга на коротком горизонте.

\begin{figure}[t]
\centering
\includegraphics[width=\linewidth]{fig5_social_channel.pdf}
\caption{Влияние экономических условий на социальное напряжение. \emph{(a)}~Импульсный отклик среднего
напряжения на экзогенный стагфляционный шок (безработица и инфляция, три года): средняя кривая и
полоса $5$--$95$-го перцентиля по приорам; отклик нарастает за $\approx4$~года и сохраняется.
\emph{(b)}~Локализация: ранговая корреляция приоров с итоговым напряжением; доминируют параметры
социального блока (помечены $*$), а не физико-экономические.}
\label{fig:social}
\end{figure}

\section{Ограничения}
\label{sec:limits}

Откровенный перечень того, чего модель пока \emph{не} может, не менее важен, чем перечень
её возможностей.

\textbf{Ожидания в базовом прогоне адаптивны.} Домохозяйства и фирмы ориентируются на
недавнее прошлое, а не на рационально ожидаемое будущее, поэтому базовый прогон не воспроизводит
эффекты опережающего поведения (например, реакцию на заранее объявленную политику до её
вступления в силу). Модельно-состоятельный (near-rational) механизм ожиданий реализован и
протестирован, однако ансамблевая проверка показала, что его маржинальный эффект статистически
неустойчив, поэтому он оставлен отключённым по умолчанию.

\textbf{Долгосрочный рост слабо идентифицирован.} Темп долгосрочного роста производительности
задаётся эндогенным механизмом (накопление НИОКР по Джонсу, переток технологий, сближение с
границей), но лишь слабо определяется имеющейся короткой историей и для прогнозов вперёд привязан
к внешнему сценарию, а не выводится из неё.

\textbf{Климат не описан фундаментальной моделью.} Климатический блок воспроизводит общепринятые
показатели, но не является полной моделью земной системы; сильнейшие углеродные обратные
связи глубоко неопределённы и вынесены в хвостовые сценарии. Тем не менее это представляется рабочей аппроксимацией для стратегирования и управления рисками.

\textbf{Оценки ущерба неопределённы во всей дисциплине.} Оценки климатического ущерба
неопределённы во всём поле исследований; модель несёт эту неопределённость как интервал, а
не как единое число.

\textbf{Социально-политические связи частично заякорены, частично экспертны.} Единого
калибровочного набора для перехода от условий к недовольству не существует, однако значительная
часть приоров социального блока теперь привязана к литературе и воспроизводимым числам
(раздел~\ref{sec:uq}): показатель степенного хвоста размеров войн, эластичности гравитации миграции,
величина обвала ВВП при коллапсе режима и прогностическая надёжность начала конфликта, --- а
географическое сопряжение проверено против пространственной зависимости реального реестра.
Остаточно-экспертными остаются пороги связей и, в частности, вес приграничной диффузии напряжённости:
чистого приграничного якоря для неё нет --- в наиболее аккуратном межстрановом исследовании
приграничный перелив протестов статистически незначим. Сам канал влияния экономики на напряжение
значим, направленно верен и запаздывает (рис.~\ref{fig:social}); конфликт же по своей природе плохо
предсказуем, и геополитический блок пригоден лишь для относительного риска (AUC ${\approx}0{,}74$).

\textbf{Триггеры кризисов основаны на порогах.} Дискретные кризисы запускаются пороговыми правилами
и схватывают верное качественное поведение, не будучи выведены из единого подогнанного набора
данных о кризисах; распределение их тяжести при этом привязано к степенному закону размеров войн
(показатель ${\alpha}{\approx}1{,}5$; Ричардсон; Клозет).

\textbf{Географическое сопряжение реалистично, но выигрыш сосредоточен в торговле.} Распространение
шоков заземлено на реальный пространственный граф, и модель воспроизводит пространственную
зависимость реального реестра (надбавка конфликтности у соседей, эмерджентный индекс Морана). Однако
содержательный выигрыш сосредоточен в торговом канале (заякоренная гравитация), тогда как веса
приграничной диффузии конфликта, напряжённости и климатического риска скромны и для напряжённости
лишены чистого литературного якоря.

Общая установка остаётся неизменной: реалистичная цель --- быть откалиброванным
относительно собственной неопределённости, а не служить инструментом безошибочного
предвидения. Реалистичной планкой точности многолетних прогнозов служит превосходство над
наивным эталоном.

\section{Типовые сценарии применения}
\label{sec:scenarios}

Модель предназначена для трёх взаимосвязанных режимов работы, в которых её сопряжённость и
корректная оценка неопределённости дают наибольшую отдачу.

\textbf{Сценарный анализ и стратегическое планирование.} Основной режим --- сравнение базовой
траектории со сценарием возмущения или решения политики: вводится шок или мера (например, цена
на углерод, ресурсный дефицит, обострение отношений между странами), и прослеживается её
распространение по экономике, климату, обществу и геополитике. Это позволяет сопоставлять
альтернативные стратегии и оценивать сопутствующие риски. Содержательно значимы здесь
\emph{относительные} величины --- разности между сценарием и базой, --- а не абсолютные
точечные значения на дальнем горизонте.

\textbf{Ансамблевые расчёты.} Поскольку расчётное ядро детерминировано, неопределённость
вводится прогоном больших ансамблей со случайным выбором параметров из обоснованных
диапазонов. Это даёт не одну траекторию, а распределение исходов, по которому измеряется
чувствительность к каждому параметру и строятся вероятностные оценки. Сценарии
рекомендуется сравнивать именно как ансамбли, а не как одиночные прогоны.

\textbf{Обнаружение слабых сигналов.} Встроенная логика кризисов имеет пороговый характер: кризис
срабатывает при пересечении границы. Обнаружение слабых сигналов --- дополнение к ней:
выявление малозаметного многомерного накопления напряжения, \emph{предшествующего}
структурному сдвигу, до того как сработает любой порог. Для этого применяются три
взаимодополняющих средства. Первое --- расстояние Махаланобиса: для каждого момента
времени оценивается, насколько \emph{совместное} состояние (одновременно продукт,
температура, долг, напряжение и т.\,д.) удалено от его недавнего базового распределения с
учётом корреляций между измерениями; превышение порога $\chi^2$ указывает на переход
системы в необычную область пространства состояний \citep{mahalanobis1936,wilson1931}.
Второе --- обнаружение структурного сдвига: на отдельном ряду ищется смена режима методом
обнаружения разладки со штрафом по информационному критерию, что даёт вероятность сдвига и
его наиболее вероятное положение \citep{schwarz1978}. Третье --- индикаторы критического
замедления (рост автокорреляции и дисперсии), нарастающие при приближении к точке
бифуркации \citep{scheffer2009}. Чтобы эти средства выявляли отклонения от \emph{динамики},
а не вторично детектировали плавный тренд, ряды предварительно приводятся к стационарному
приращению. На синтетических данных и на траекториях самой модели средства ведут себя
ожидаемо: внесённый спад продукта на $15\,\%$ распознаётся точно на нужном шаге, тогда как
спокойный прогон остаётся вблизи теоретической частоты ложных срабатываний. Рабочий приём
здесь, как и в сценарном анализе, --- \emph{относительный}: сопоставление профиля слабых
сигналов сценария с базовым (более ранние и более частые отклонения, более высокая
вероятность сдвига) указывает, что сценарий смещает систему к структурному переходу.

\section{Программная реализация и сценарная среда}
\label{sec:repo}

Модель реализована как открытая кодовая база и доступна в репозитории
\url{https://github.com/Theclimateguy/GIM}. Расчётное ядро детерминировано и
воспроизводимо: каждый прогон записывается в отдельную датированную папку с манифестом
параметров и версий входных данных, а набор модульных тестов и протокол валидации фиксируют
ключевые величины ретроспективной проверки (раздел~\ref{sec:calib}). Калибровочные наборы
данных подключаются как неизменяемые входы, привязанные к контрольной сумме, что исключает
незаметное расхождение результатов между запусками. Базовый мировой прогон, сценарные
вопросы, многоагентные партии и калибровка запускаются единым интерфейсом командной строки.

\textbf{Большие языковые модели в роли стран.} Помимо штатных правил поведения, политику
отдельной страны может задавать агент на основе большой языковой модели. Агент получает
наблюдаемое состояние своей страны и её окружения и возвращает связную стратегию ---
внутреннюю, внешнюю, финансовую, военную и торговую политику, --- которая компилируется в
формализованные действия модели. Поддерживаются режимы фоновой политики: правило, целевой
рост, языковая модель и предварительно скомпилированная доктрина; глубину обновления
стратегии (по событию, периодически или однократно) задаёт пользователь. Здесь важно
методологическое разграничение: генеративная модель порождает только намерения и стратегию
акторов, тогда как их последствия рассчитывает детерминированное, эмпирически закреплённое
ядро. Языковая модель отвечает за правдоподобие поведения, имитационное ядро --- за
правдоподобие механики и числовых исходов; такое разделение сохраняет воспроизводимость и
ограждает количественные результаты от произвольных суждений языковой модели.

\textbf{Гибридные партии и гипотезы «что, если».} Среда допускает совместную игру человека и
агентов: часть стран ведут пользователи, часть --- языковые агенты или штатные правила, а
намерения задаются на естественном языке и переводятся в действия модели. Предусмотрены два
режима хода. В режиме \emph{действия} ход фиксируется и переносится в дальнейшую траекторию;
в режиме \emph{«что, если»} гипотеза проигрывается обособленно, без необратимых последствий
для основной линии, что удобно для проверки контрфактических предположений (например, «как
изменится обстановка, если страна~X ужесточит санкции против страны~Y»). Каждый ход
разворачивается не в одну траекторию, а в ансамбль, так что отклик измеряется вместе с его
неопределённостью --- как и в сценарном анализе раздела~\ref{sec:scenarios}.

\textbf{Интерактивная аналитическая среда.} Отдельная команда поднимает локальную
аналитическую панель, в которой собраны сценарные вопросы, многоагентные партии и метрики
стран; вкладка партий использует описанный гибридный механизм. Результаты каждого прогона
выгружаются в наглядном виде --- сводная панель, аналитическая записка для лиц, принимающих
решения, и журналы политик и базовых траекторий, --- что делает многоагентный эпизод
прослеживаемым и пригодным для последующего разбора.

Тем самым сценарная среда операционализирует заявленное назначение модели. Сопряжённое ядро
превращается из расчётного механизма в площадку для отработки многоагентного взаимодействия:
коллектив языковых агентов или смешанный состав «человек--агент» проигрывает конкретные
сценарии и контрфактические гипотезы, а пользователь наблюдает распространение шоков по
экономике, климату, обществу и геополитике и ранние слабые сигналы приближающихся
структурных сдвигов --- в едином, воспроизводимом и поддающемся аудиту контуре. Роль
генеративного ИИ здесь вспомогательная, но содержательная: он расширяет пространство
правдоподобных стратегий акторов, не подменяя собой эмпирически закреплённую механику
последствий.

\section{Заключение}

Представленная модель связывает в едином контуре экономику, климат, природный капитал,
общество, политику, геополитику и культуру и тем самым заполняет содержательный пробел
между специализированными моделями отдельных областей. Её ценность для стратегического
планирования, сценарного анализа и управления рисками состоит не в точечном прогнозе, а в
способности связно проигрывать сценарии, сопоставлять альтернативные решения, корректно
выражать неопределённость через ансамбли и выявлять слабые сигналы приближающихся
структурных сдвигов. Экономический и климатический блоки эмпирически закреплены и проходят
ретроспективную проверку; социально-политические и геополитические блоки, составляющие
отличительную черту модели, проверены по более скромному, но содержательному критерию ---
превосходству над наивными эталонами. Экономический контур с тех пор углублён: добавлены явная передача от денежной массы к ценам
(откалиброванный член количественной теории), эндогенный механизм долгосрочного роста (накопление
НИОКР по Джонсу) и протестированный механизм опережающих ожиданий, оставленный отключённым, так как
его маржинальный эффект оказался неустойчивым. Остаточные ограничения --- приближённость
климатического блока, дисциплинарная неопределённость оценок ущерба и принципиально низкая
предсказуемость конфликта --- изложены честно и не препятствуют применению модели в очерченных режимах.

\newpage

\appendix
\section{Ключевые параметры и их идентификация}
\label{app:params}

Каждый параметр задан априорным распределением; ниже приведены центральное значение,
диапазон правдоподобия (примерно $5$--$95$-й перцентиль приора) и основной источник. $33$~ключевых параметра
снабжены литературно-обоснованными приорами (у каждого --- источник и обоснование); полный реестр
модели насчитывает $287$~параметров с тегами статуса калибровки, большинство из которых заданы
экспертно или структурно, --- он доступен в репозитории (раздел~\ref{sec:repo}).

\begin{table}[h]
\centering
\small
\begin{tabularx}{\linewidth}{@{}lccL@{}}
\toprule
Параметр & Центр. & Диапазон & Источник \\
\midrule
Равновесная чувствительность климата, $^\circ$C & $3{,}0$ & $1{,}5$--$6{,}0$ & МГЭИК AR6 WG1 (2021); Otto и др.\ 2013 \\
Теплоёмкость поверхности & $8{,}0$ & $6{,}0$--$18{,}0$ & Geoffroy и др.\ 2013 (двухслойная EBM) \\
Обмен теплом с океаном & $1{,}0$ & $0{,}4$--$1{,}3$ & Geoffroy и др.\ 2013 \\
Масштаб эмиссий & $0{,}98$ & $0{,}90$--$1{,}05$ & Global Carbon Project 2023 \\
Темп структурной декарбонизации & $0{,}016$ & $0{,}010$--$0{,}030$ & Тренд углеродоёмкости GCP\,/\,WDI 2000--2023; бэктест GIM 2015--2023 \\
Доля капитала $\alpha$ & $0{,}30$ & $0{,}22$--$0{,}40$ & Penn World Table 10; Gollin 2002 \\
Доля труда $\beta$ & $0{,}60$ & $0{,}50$--$0{,}72$ & Penn World Table 10; Gollin 2002 \\
Эластичность по энергии $\gamma$ & $0{,}042$ & $0{,}015$--$0{,}080$ & Бэктест 2015--2023; лит-ра по энергии \\
Замещение капитал--энергия $\sigma_{KE}$ & $0{,}40$ & $0{,}20$--$0{,}60$ & van der Werf 2008; Koetse и др.\ 2008 \\
Норма амортизации $\delta$ & $0{,}05$ & $0{,}03$--$0{,}08$ & Penn World Table 8/9 \\
Чувствительность ПФП к НИОКР & $0{,}30$ & $0{,}10$--$0{,}60$ & Калибровка по бэктесту \\
Коэффициент ущерба & $0{,}006$ & $0{,}0015$--$0{,}025$ & DICE-2016R2; Говард и Стернер 2017; Бёрк и др.\ 2015 \\
Норма временного предпочтения $\rho$ & $0{,}015$ & $0{,}001$--$0{,}030$ & DICE-2016R2; Stern 2007 \\
Эластичность предельной полезности $\eta$ & $1{,}45$ & $1{,}0$--$2{,}0$ & DICE-2016R2; Stern 2007 \\
Эластичность рыночного спроса & $0{,}40$ & $0{,}20$--$0{,}70$ & Labandeira и др.\ 2017 (мета-анализ) \\
\midrule
\midrule
\multicolumn{4}{@{}l}{\emph{Социальный блок (структурные допущения; проверяются по знаку и каналу, раздел~\ref{sec:uq})}}\\
Стресс напряжения от безработицы & $0{,}010$ & $0{,}007$--$0{,}014$ & Структурное допущение (рост безработицы повышает недовольство) \\
Стресс напряжения от инфляции & $0{,}005$ & $0{,}003$--$0{,}007$ & Структурное допущение \\
Эффект неравенства на напряжение & $0{,}0005$ & $0{,}0003$--$0{,}0007$ & Структурное допущение \\
Сила якоря доверия & $0{,}060$ & $0{,}039$--$0{,}081$ & Структурное допущение (возврат к равновесию) \\
\bottomrule
\end{tabularx}
\end{table}

\clearpage
\section{Бенчмарк интеграции: где останавливаются секторальные модели}
\label{app:benchmark}

Утверждение о широте сопряжения (таблица~\ref{tab:compare}) здесь сделано вычислимым. Схема
сравнения опирается на литературу, а не на внутреннее отключение блоков: для каждого из трёх
шоков берётся опубликованный вывод конкретной секторальной модели, той же модели GIM подаётся
тот же экзогенный шок, и проверяется, во-первых, воспроизводит ли GIM первопорядковый эффект
секторальной модели и, во-вторых, обнаруживает ли она межсекторное последствие, которое
секторальная модель не способна представить по самому своему устройству. Предметом сравнения
служит вывод, значимый для решения, а не отдельное число: не утверждается, что GIM точнее DICE
считает выбросы или благосостояние, --- утверждается, что она видит последствие, к которому DICE
структурно слепа. Прогоны детерминированы (полный мир из пятидесяти семи стран, горизонт десять
лет, экстремальные события отключены, так что единственное различие между базовым и шоковым
прогонами --- введённый шок). Каждый сценарий сопровождается зависимостью отклика от величины
шока, на которой наклон секторальной модели по построению равен нулю. Расчётный стенд и рисунки
воспроизводимы из репозитория (раздел~\ref{sec:repo}).

Два методологических замечания приводятся прямо, поскольку они входят в результат. Во-первых,
агрегированная подсистема цен на ресурсы насыщается уже в пределах первого расчётного года,
поэтому чисто ценовой шок не распространяется вперёд; каждый шок вводится в ту переменную, через
которую он впервые выходит за пределы своего сектора, --- именно в тот канал, которого нет у
секторальной модели: для углеродного налога это перенос налога в потребительские цены, для
нефтяного шока --- бремя импортного счёта и расход резервов, для урожайного шока --- разрыв в
продовольственном обеспечении. Во-вторых, три перетока различаются по величине и форме отклика и
сообщаются как есть, без выравнивания в большую сторону.

\subsection{Углеродный налог (эталон: DICE)}

Модель DICE-2016R \citep{nordhaus2017} находит оптимальную по благосостоянию цену углерода (около
тридцати одного доллара за тонну в 2015 году с последующим ростом), единственный первопорядковый
эффект которой --- снижение выбросов; в ней нет ни рынка труда, ни политической экономии. При том
же налоге в пятьдесят долларов за тонну GIM воспроизводит снижение выбросов (около пяти с половиной
процентов за десять лет) и дополнительно повышает социальное напряжение и снижает доверие к власти: перенос
налога в цены на энергию питает инфляцию и безработицу. Это тот политэкономический механизм,
который стоит за отменой французского углеродного и топливного налога в 2018 году и австралийской
цены на углерод в 2014 году. Это канал, которого у DICE нет по построению: по всей развёртке ставки
отклик напряжения остаётся ненулевым (рис.~\ref{fig:bench-carbon}) --- малым и немонотонным в этом
масштабе, но отличным от тождественно нулевого отклика DICE.

\begin{figure}[t]
\centering
\includegraphics[width=\linewidth]{fig6_integration_carbon.pdf}
\caption{Сценарий углеродного налога. Слева: GIM воспроизводит снижение выбросов, которое моделирует
DICE. В центре: GIM добавляет рост социального напряжения, недоступный DICE. Справа: зависимость
отклика напряжения от ставки налога ненулевая у GIM (хотя малая и немонотонная) против нулевого отклика у DICE.}
\label{fig:bench-carbon}
\end{figure}

\subsection{Нефтяной ценовой шок (эталон: энергетические модели)}

Энергосистемная модель, такая как MESSAGEix, разрешает шок предложения нефти как рост цены и
подстройку спроса в пределах энергетического сектора. GIM проводит тот же шок дальше, в суверенные
финансы: возросший счёт за импорт, покрываемый расходованием резервов и наращиванием долга,
подводит наиболее уязвимых в бюджетном отношении импортёров к порогам долгового кризиса, причём
какие именно страны срываются, определяют собственные пороги модели по долгу и процентной ставке,
а не ручной отбор. При реалистичном шоке (сокращение предложения на двадцать процентов, бремя
порядка трёх процентов ВВП) эффект краевой --- в кризис попадает один дополнительный импортёр; при
тяжёлом шоке (бремя порядка пятнадцати процентов ВВП, сопоставимое с худшими историческими
нефтяными эпизодами) средняя длительность долгового кризиса у импортёров возрастает примерно в два
с половиной раза, а в суверенный кризис входят ещё четыре страны (рис.~\ref{fig:bench-oil}).
Отклик имеет пороговый, ступенчатый характер, а энергосистемная модель не порождает суверенных
кризисов ни при какой величине шока. Этот механизм стоит за эпизодами вроде дефолта Шри-Ланки в
2022 году.

\begin{figure}[t]
\centering
\includegraphics[width=\linewidth]{fig7_integration_oil.pdf}
\caption{Сценарий нефтяного ценового шока. Слева: длительность суверенного долгового кризиса у
импортёров при тяжёлом шоке. Справа: зависимость отклика от величины импортного бремени имеет
пороговый, ступенчатый характер; отмечены краевая (реалистичная) и выраженная (тяжёлая) точки против
нулевого отклика энергосистемной модели.}
\label{fig:bench-oil}
\end{figure}

\subsection{Урожайный шок (эталон: AgMIP и AR6)}

Оценки урожайности в проектах AgMIP и в шестом докладе МГЭИК останавливаются на снижении урожая и
индексе продовольственной безопасности. GIM проводит ту же потерю урожая дальше, во внутреннюю
нестабильность: сокращение производства и запасов у наиболее бедных импортёров открывает разрыв в
продовольственном обеспечении, который повышает протестное давление, причём по мере углубления
потерь эскалация ускоряется. Это та нелинейная «граница высоких последствий», что задокументирована
для цен на продовольствие и «Арабской весны» \citep{lagi2011food} (порог индекса продовольственных
цен ФАО на отметке двести десять). При тяжёлой потере урожая в пятьдесят процентов на уязвимом
подмножестве стран стресс продовольственной доступности возрастает примерно на девять сотых, а
протестное давление --- примерно на тринадцать тысячных (рис.~\ref{fig:bench-crop}); далее по цепочке
отклик затухает к хрупкости режима, и конечной точкой оказывается внутренняя нестабильность, а не
межгосударственный конфликт.

\begin{figure}[t]
\centering
\includegraphics[width=\linewidth]{fig8_integration_crop.pdf}
\caption{Сценарий урожайного шока. Слева: стресс продовольственной доступности и протестное давление
на уязвимом подмножестве стран. Справа: зависимость отклика от глубины потери урожая выпукла
(ускоряется) против нулевого отклика AgMIP по оси протестного давления.}
\label{fig:bench-crop}
\end{figure}

Итог. По всем трём каналам (рис.~\ref{fig:bench-summary}) GIM демонстрирует положительный измеримый
наклон отклика по той оси, на которой любая секторальная модель по построению остаётся на нуле.
Именно этот зазор --- а не отдельное число --- и составляет выигрыш от интеграции; это ровно то
измерение, которое таблица охвата описывает качественно.

\begin{figure}[t]
\centering
\includegraphics[width=\linewidth]{fig9_integration_summary.pdf}
\caption{Выигрыш от интеграции. Каждая панель --- один межсекторный канал, нормированный к
собственному максимуму; пунктир --- нулевой отклик любой секторальной модели по этой оси. Наклон GIM
плавный для канала «экономика и общество», ступенчато-пороговый для канала «ресурсы и финансы» и
выпуклый для канала «климат и общество».}
\label{fig:bench-summary}
\end{figure}

\newpage
\bibliographystyle{plainnat}
\bibliography{references}

\end{document}
